% What @interfaceObject does with a type name. Referenced from chapter 13. Bare
% tikzpicture; the figure environment, caption and label live at the call site.
%
% The point of the picture is the label on the middle arrow. The router reads
% keys off concrete types it knows about and sends them typed as the interface,
% so the subgraph on the right answers about a book without ever being told it
% is answering about a book.
\begin{tikzpicture}[
    font=\small,
    svc/.style={draw, rounded corners=2pt, minimum width=26mm,
                minimum height=15mm, align=center, font=\scriptsize},
    tp/.style={draw, rounded corners=2pt, minimum width=15mm,
               minimum height=5mm, align=center, font=\scriptsize},
    flow/.style={-{Stealth[length=2mm]}},
    note/.style={font=\scriptsize, gray, align=center, inner sep=1pt}
  ]

  % -- left: the owner --------------------------------------------------------
  \node[note, anchor=south] at (1.7,1.5) {\code{library}};
  \draw[rounded corners=2pt] (0.2,-1.9) rectangle (3.2,1.4);

  \node[tp] (media) at (1.7,0.9) {\code{interface Media}};
  \node[tp] (book)  at (1.7,-0.4) {\code{Book}};
  \node[tp] (film)  at (1.7,-1.3) {\code{Film}};

  \node[note, anchor=north, align=center] at (1.7,-2.1) {both entities,\\keyed on \code{id}};

  % -- right: the contributor -------------------------------------------------
  \node[note, anchor=south] at (8.6,1.5) {\code{ratings}};
  \draw[rounded corners=2pt] (7.1,-1.9) rectangle (10.1,1.4);

  \node[tp, minimum width=22mm] (io) at (8.6,0.5) {\code{type Media}\\\code{@interfaceObject}};
  \node[note, anchor=north, align=center] at (8.6,-0.8)
    {no \code{Book} here,\\no \code{Film},\\no way to list them};

  % -- the wire ---------------------------------------------------------------
  \draw[flow] (3.3,0.2) -- node[note, above, align=center]
    {keys read off\\\code{Book} and \code{Film}} node[note, below, align=center]
    {sent as\\\code{\_\_typename: Media}} (7.0,0.2);

\end{tikzpicture}
